\subsection{风格训练原理与模型设计}
\subsubsection{从风格统计到前馈网络}
神经风格迁移同时约束内容结构和绘画统计。内容约束使输出保持人物、物体及空间关系；风格约束使特征相关性接近参考画作。与对每张图反复优化像素的方法相比，前馈网络在训练阶段学习这一映射，推理阶段只执行固定算子\cite{gatys,johnson}。本系统的风格参考只在训练时参与损失计算；FPGA推理输入始终只有RGB三通道，风格通过选择独立权重确定。

\reportfigure{assets/generated/training.png}{1}{训练、量化与部署流程示意。VGG和教师只用于离线训练，板端执行六通道学生网络。}{fig:training}
令内容图为$x$，第$s$种风格输出为$y=f_{\theta_s}(x)$，教师输出为$t_s(x)$。冻结VGG的第$l$层特征为$\phi_l$，Gram统计定义为
\begin{equation}
 G_l(z)=\frac{\Phi_l(z)\Phi_l(z)^{\mathsf T}}{C_lH_lW_l},\qquad
 \Phi_l(z)=\operatorname{reshape}\bigl(\phi_l(z)\bigr).
\end{equation}
内容项比较输出与输入的高层特征；风格项比较输出与参考画作$r_s$在若干层的Gram矩阵：
\begin{equation}
 L_c=\|\phi_c(y)-\phi_c(x)\|_2^2,\qquad
 L_s=\sum_l\|G_l(y)-G_l(r_s)\|_2^2.
\end{equation}
实现使用各张量的平均误差，不将大分辨率的特征元素总数直接作为损失尺度。

仅强化逐像素拟合容易把绘画教师压成颜色滤镜。本轮降低像素蒸馏约束，并比较局部平均梯度幅值，使学生保留细笔触和局部纹理。令$E_k$分别为对水平、垂直绝对梯度作8×8局部平均的算子，$k\in\{x,y\}$，
\begin{equation}
 L_e=\sum_k\|\log(E_k(y)+\epsilon)-\log(E_k(t_s)+\epsilon)\|_1 .
\end{equation}
训练还使用像素蒸馏$L_d=\|y-t_s\|_1$、8倍平均池化后的粗结构项$L_p$和TV平滑项。当前微调代码的损失组合为
\begin{equation}
 L=0.8L_d+L_p+0.12L_c+220000L_s+0.25L_e+0.025L_{\rm TV}.
\end{equation}
水墨槽额外加入通道离均值的色度惩罚，系数为2，以抑制彩色背景。该系数对应当前Gram归一化和特征层选择，不能离开实现直接与其他论文的系数大小比较。


\blocktitle{各项损失解决什么问题}
\begin{table}[H]\centering\small
\caption{当前训练目标、输入与视觉作用}
\begin{tabularx}{\linewidth}{@{}p{.18\linewidth}p{.30\linewidth}X@{}}\toprule
项 & 比较对象 & 作用及过强时的风险\\\midrule
像素蒸馏 & 学生与对应教师RGB & 锚定构图和基本色彩；过强易平均掉笔触\\
粗结构 & 8倍平均池化学生/教师 & 保留大块明暗，不逐像素强行对齐纹理\\
VGG内容 & 学生与原图第3组特征 & 约束主体结构；过强趋近原图\\
Gram风格 & 学生与参考图4组统计 & 学习颜色与特征共现；不提供物体语义布局\\
局部梯度 & 学生与教师8×8梯度均值 & 鼓励短笔触；过强可能把背景推成噪纹\\
TV & 输出相邻像素差 & 限制高频振荡；过强又会抹平磨砂感\\
水墨色度 & RGB对通道均值偏差 & 压低彩色残留；不等价于重新设计墨色结构\\\bottomrule
\end{tabularx}
\end{table}
VGG取冻结特征序列索引3、8、17、20，内容使用第3组；进入VGG前按ImageNet均值和标准差归一化。Gram矩阵省去空间位置，能使局部统计接近参考，却不要求星空旋涡位于天空或人物之外。因此“使用星月夜参考”不等于生成具有同等构图的大型旋涡，这既与损失定义相关，也受RF31限制。

局部纹理项以$\epsilon=0.012$稳定对数，比较水平、垂直绝对梯度的8×8平均，并非多尺度频域纹理损失。对数使弱纹理也有贡献，降低仅追逐最大边缘的倾向。该项配合较弱像素回归用于修复“只有蓝色滤镜”的退化，但当前没有做足够的同预算消融实验，不能将画质改善全部归因于单一损失。

\subsubsection{全分辨率六通道结构}
\reportfigure{assets/generated/network.png}{1}{学生网络示意。图中省略各卷积后的冻结BN与定点激活；输入和隐藏层的空间尺寸保持不变。}{fig:network}
输入投影为1×1卷积，RGB三通道转为六通道。四个残差块分别执行DW 3×3和PW 1×1，膨胀率为1、2、4、8；最终有符号1×1输出头产生RGB增量，与原始RGB相加后饱和。深度可分离卷积的空间与通道分解降低运算量\cite{mobilenet}，本模型未照搬完整MobileNet分类网络。

\begin{table}[H]\centering\small
\caption{单风格网络结构及卷积权重数量}
\begin{tabularx}{\linewidth}{@{}p{.27\linewidth}Xrr@{}}\toprule 阶段 & 算子与通道 & 膨胀率 & 权重数\\\midrule
输入投影 & 1×1，3→6 & 1 & 18\\
残差块1 & DW 3×3六通道＋PW 1×1 6→6 & 1 & 90\\
残差块2 & 同上 & 2 & 90\\
残差块3 & 同上 & 4 & 90\\
残差块4 & 同上 & 8 & 90\\
输出头 & 1×1，6→3，有符号 & 1 & 18\\\midrule
合计 & 十个卷积算子 & — & 396\\\bottomrule
\end{tabularx}
\end{table}
同为六输入、六输出通道时，普通3×3卷积每像素需要$9C^2=324$次乘加，DW加PW需要$9C+C^2=90$次，减少约72.2\%。这只是同尺寸同通道的算子比较；残差缓存、数据搬运和并行度仍影响整机资源。

DW权重为$9C$、PW为$C^2$；$C=6$时每残差块为54+36=90。感受野为$1+2(1+2+4+8)=31$。训练模型单分支可训练参数共567个，含BN、仿射及偏置；396专指卷积系数，不能当成所有训练参数。十分支共5670个可训练参数。部署折叠后每风格保存396个INT8系数、57组偏置与移位及对齐填充，有效712B、包跨度736B，十组共7360B。只有参数分十组，硬件计算核共用。

训练的144×144随机裁剪用于批训练，并不把推理降为144×144。模型为全卷积，部署与本轮对照均在1280×720运行，无池化、2×2共享输出或固定XY特征。条带调度是实现层的存储安排，不改变全图卷积定义。

\subsubsection{训练数据与本轮微调}
训练固定512张内容图，24张验证图，与34张用户留出照片通过指纹排除，训练和验证路径不重叠。四个风格使用保留的绘画教师与参考，其他六类以既有C8整数教师作较弱空间锚点，并接受对应风格参考的VGG约束；不是十个风格全部从高容量绘画教师训练。

本轮先对梵高试训1500步，再进行十槽12000步纹理微调，batch为4、内容裁剪144×144、参考图384×384，Adam初始学习率$5.5\times10^{-5}$。12000步调度中梵高更新3693次，其余各923次；加上试训，梵高新增5193次。本轮新增总计13500步，checkpoint的全局step27500包含继承的历史计数，不表示每槽都训练27500次。图像增强包括随机裁剪、水平翻转、曝光及gamma变化，并量化输入至RGB565。

教师参考中，部分剪纸、漫画图为此前选定的AI概念素材；风格名称用于标识参考方向，不把概念图说成名家原作。本报告的两张公开名画只用于训练后评估展示，不参与本轮微调。

\reportfigure{assets/generated/training_log.pdf}{1}{真实12000步微调日志。左图为选定四槽40次更新滑动平均训练损失，右图为各槽更新数；不是验证损失或画质评分。}{fig:training-log}


\blocktitle{一次训练迭代的可复现流程}
每次从固定内容缓存抽取4张，经随机裁剪、翻转、亮度与gamma处理，再按RGB565量化并扩回RGB。该输入同时送教师与学生，避免教师使用高精度颜色、学生使用板端低精度颜色造成域差异。内容图先整理为176×176缓存，迭代取144×144；风格参考单独缩放到384×384计算并缓存Gram，不与内容拼成六通道输入。

十槽调度序列为$[0,1,0,2,3,0,4,5,6,0,7,8,9]$，梵高在每13步出现4次，其他槽各1次。一次仅对被选分支反向传播，教师和VGG不更新；总损失求导后将梯度范数限制为4，再进行Adam更新和系数范围投影。QAT卷积路径关闭混合精度，教师/VGG可用bfloat16降低训练开销，不能把VGG训练精度直接当成FPGA算术格式。

十组参数从完成量化感知训练的学生网络导出。训练损失曲线只能说明优化过程，不代表主观画质、泛化性能或实际帧率。

\subsubsection{冻结BN与整数QAT}
冻结BN的统计量折叠进卷积：
\begin{equation}
 \bar W_c=\frac{\gamma_c}{\sqrt{\sigma_c^2+\varepsilon}}W_c,\qquad
 \bar b_c=\beta_c-\frac{\gamma_c\mu_c}{\sqrt{\sigma_c^2+\varepsilon}}.
\end{equation}
若有学习仿射，再将其增益和偏置合并。它是冻结BN的确定性部署，不等同于在硬件实时计算IN。

逐输出通道选择二进制移位$s_c$，将实系数缩放、按$\lfloor z+0.5\rfloor$舍入并量化到$[-127,127]$。uint8输入与有符号INT8权重相乘，偏置与累加使用有符号算术：
\begin{equation}
 a_c=\sum_i w_{ci}x_i+b_c,\quad
 q_c=\left\lfloor\frac{a_c+2^{s_c-1}}{2^{s_c}}\right\rfloor\;(s_c>0).
\end{equation}
$s_c=0$时不加半量直接取$a_c$。负数也遵循该式，不改成绝对值对称舍入。中间激活饱和到$[0,255]$，头部为$[-510,510]$，最后与原RGB相加并饱和到$[0,255]$。残差支路同样按量化尺度对齐。

QAT前向使用与硬件一致的折叠、量化、舍入、饱和，反向用直通估计器近似取整梯度\cite{qat}。参数导出同时计算
\begin{equation}
 B_c=|b_c|+255\sum_i|w_{ci}|
\end{equation}
约束累加位宽。当前最大保守界239482，小于有符号22位上限；加舍入量的通路采用23位。实际乘积、加法树、残差和输出逐层与Python整数参考核对，不仅比较最终视觉效果。


\blocktitle{真实量纲与整数编码}
输入RGB以$S_{\rm rgb}=1/255$编码，隐藏特征以$S_h=6/255$编码。设实数输入$x_{\rm real}=S_{\rm in}x_q$，期望输出$y_{\rm real}=S_{\rm out}y_q$，则量化前系数应为
\[
c_{oi}=\bar W_{oi}\frac{S_{\rm in}}{S_{\rm out}},\qquad
d_o=\frac{\bar b_o}{S_{\rm out}}.
\]
stem使用$S_{\rm in}/S_{\rm out}=1/6$；DW/PW中间层为1；输出头为6。先做这个尺度变换，再选通道移位并量化，不能把同一浮点权重直接乘255用于所有层。各块残差两支均在$6/255$尺度，故可直接整数相加再截断；最终RGB旁路和头部增量都在$1/255$尺度。

每通道移位初值由最大系数决定，约为$\lfloor\log_2(127/\max_i|c_{oi}|)\rfloor$，限制0--20，并按累加上界进一步降低。较大移位可保留小权重精度，但会放大偏置与累加数值，必须同时考虑位宽；只做INT8系数截断会忽略偏置溢出风险。当前导出器对十槽逐层检查，系数裁剪计数为0；这不是所有未来模型都天然满足的性质。

例如某输出通道$s=3$，累加$a=-12$，硬件取$(-12+4)\gg3=-1$；当$a=12$时取$(12+4)\gg3=2$。这说明它采用加半量后算术右移，与“四舍五入绝对值后恢复符号”不同。定点参考必须复现符号扩展、舍入和饱和顺序，否则边缘纹理会出现大量一位差异。

\subsubsection{从训练产物到FPGA}
\begin{enumerate}
\item 冻结最终checkpoint，逐槽折叠、量化并导出权重、偏置和移位；重排成每层的配置事务。
\item 生成736B对齐包和只读ROM；RISC-V按层、类型、索引发出216项配置，确认后commit。
\item 先对照软件整数与QAT，再检查RTL、背压、条带边界、实际映射算子及固件加载。
\item 通过整机综合、布局布线、时序分析与引脚检查，生成包含模型参数和控制固件的配置文件。
\end{enumerate}
当前权重“运行时加载”是CPU从预置ROM装入加速器，不是从电脑或网络下载新权重。外部在线更新作为后续目标在第四部分单独说明。

\blocktitle{712B数据包与216次写入如何得到}
每层包依次存INT8权重、补齐到4B边界的零字节、小端INT32偏置、UINT8移位及其4B对齐。包对齐与总线写入数不同：偏置/移位逐通道各发一次32位配置；ROM内移位仍只占1B。
\begin{table}[H]\centering\small
\caption{单风格参数包和配置写入构成}
\begin{tabular}{@{}lrrrrr@{}}\toprule
层组 & 数量 & 系数/层 & 通道/层 & 包字节/层 & 写次数/层\\\midrule
stem &1&18&6&52&17\\
DW &4&54&6&88&26\\
PW &4&36&6&68&21\\
head &1&18&3&36&11\\\midrule
合计 &10&396&57&712&216\\\bottomrule
\end{tabular}
\end{table}
例如DW层54个权重补齐56B，偏置24B、移位补齐8B，共88B；配置为14次权重词、6次偏置和6次移位，共26次。整模型末尾另补24B，跨度为736B。十个模型占7360B，而不是复制十套行缓存和卷积乘法器。

权重布局按输出通道、输入通道、卷积核行列顺序展平；DW每组只对应自身通道，PW核为1×1。导出器同时保留NPZ、逐层mem和打包bin，并反向拆包核对，不通过时不更新发布件。四字节词中的低字节对应较低索引系数，与CPU小端读及固定byte写者一致。

\blocktitle{数值验证与画质验证分开进行}
软件验证依次比较实际QAT前向、GPU整数运算和独立NumPy int64参考，检查每层尺寸、最终像素、边界、负累加和饱和值；RTL再覆盖输入停顿、输出背压、配置变更与条带接缝，映射原语核验证综合后的实现。当前报告新增的20个完整720p预测全部与QAT逐值一致，其中两幅梵高还进行了独立int64全帧对照。

逐值一致证明部署遵循模型，不能证明模型好看。画质检查另外关注主体可辨性、明暗层次、暗处色块、亮屏泛白、笔触细节以及切换连续性。本报告保留真实输出而不后期调色；如果两种风格趋同，应回到训练参考、损失或网络容量继续改进，而不是在展示图上补纹理。当前只给静态效果，连续视频闪烁还需要时序一致性测试。
